\section*{Problem 9}

Alphabet \(\{ \#,A,B,\ldots,Z\}\) with initial dictionary

\[
\#\mapsto 0,\quad A\mapsto 1,\ \ldots,\ Z\mapsto 26.
\]

Apply Lempel–Ziv–Welch (LZW) to the string
\[
\texttt{PETERPIPERPICKEDAPECKOFPICKLEDPEPPERS\#}.
\]

\subsection*{(a) Output codes and new dictionary entries}

Start with the initial 27 entries (codes \(0\)–\(26\)).
Scan the string, each time outputting the longest dictionary word that
matches the current input, then adding that word extended by the next
character as a new entry.

The output code sequence is:
\[
\boxed{16,\ 5,\ 20,\ 5,\ 18,\ 16,\ 9,\ 27,\ 31,\ 9,\ 3,\ 11,\ 5,\ 4,\ 1,\ 27,\ 37,\ 15,\ 6,\ 32,\ 37,\ 12,\ 39,\ 46,\ 11,\ 48,\ 4,\ 27,\ 16,\ 34,\ 19,\ 0}.
\]

New dictionary entries (beyond the initial symbols) in order are:

\[
\begin{array}{c|l}
\text{Code} & \text{Entry} \\ \hline
27 & \texttt{PE}\\
28 & \texttt{ET}\\
29 & \texttt{TE}\\
30 & \texttt{ER}\\
31 & \texttt{RP}\\
32 & \texttt{PI}\\
33 & \texttt{IP}\\
34 & \texttt{PER}\\
35 & \texttt{RPI}\\
36 & \texttt{IC}\\
37 & \texttt{CK}\\
38 & \texttt{KE}\\
39 & \texttt{ED}\\
40 & \texttt{DA}\\
41 & \texttt{AP}\\
42 & \texttt{PEC}\\
43 & \texttt{CKO}\\
44 & \texttt{OF}\\
45 & \texttt{FP}\\
46 & \texttt{PIC}\\
47 & \texttt{CKL}\\
48 & \texttt{LE}\\
49 & \texttt{EDP}\\
50 & \texttt{PICK}\\
51 & \texttt{KL}\\
52 & \texttt{LED}\\
53 & \texttt{DP}\\
54 & \texttt{PEP}\\
55 & \texttt{PP}\\
56 & \texttt{PERS}\\
57 & \texttt{S\#}
\end{array}
\]

Codes \(0\)–\(31\) fit in \(5\) bits. Once code \(32\) appears, we need \(6\)-bit codes.
The first code \(\ge 32\) in the output is \(32\), so:

- First \(19\) codes are \(<32\): \(19\times 5 = 95\) bits.
- Remaining \(13\) codes are \(\ge 32\): \(13\times 6 = 78\) bits.

Total LZW output length:
\[
\boxed{95+78=173\ \text{bits}}.
\]

\subsection*{(b) Comparison with naive encoding}

The original string has \(45\) characters.
Naively encoding each character with a fixed 5-bit code (using only the
initial alphabet) would use
\[
45\times 5 = 225\ \text{bits}.
\]

LZW uses \(173\) bits, so the savings are
\[
225-173 = \boxed{52\ \text{bits}}.
\]
